\section{\SemASE{}: Semantic Atomic Simulation Environment}
\label{sec:semase}
%% ====================================================================

To bridge the gap between simulation toolkits and the knowledge graph, we
present \SemASE{}, a Python library that extends
ASE~\cite{larsen2017ase} with RDF serialization through runtime class patching.

\subsection{Design Principles}

\SemASE{} is guided by three principles: (1)~\emph{Transparency}---minimal code
changes to existing ASE workflows; (2)~\emph{Round-trip fidelity}---objects
saved to the knowledge graph can be faithfully reconstructed; (3)~\emph{Backend
neutrality}---works with any SPARQL-compliant store.

\subsection{The \SemASE{} Ontology}

The \SemASE{} ontology maps ASE's data model to RDF:

\begin{table}[t]
\centering
\caption{Core classes in the \SemASE{} ontology.}
\label{tab:semase-classes}
\begin{tabular*}{\textwidth}{l@{\extracolsep{\fill}}ll}
\toprule
OWL Class & ASE Class & Description \\
\midrule
\texttt{semase:AtomicSystem} & \texttt{Atoms} & Collection of atoms \\
\texttt{semase:Atom} & --- & Single atom with element, position, mass \\
\texttt{semase:Cell} & \texttt{Cell} & Unit cell with lattice vectors and PBC \\
\texttt{semase:Calculator} & \texttt{Calculator} & Computational engine \\
\texttt{semase:Vector3D} & --- & 3D coordinate \\
\bottomrule
\end{tabular*}
\end{table}

The \SemASE{} ontology is designed to complement the \MLIPsOntology{}:
\texttt{semase:AtomicSystem} corresponds to individual atomic configurations
that populate a $\TrainingDatasetC$, and
\texttt{semase:Calculator} can be linked to an $\ImplementationC$.

\subsection{Usage Example}

\begin{lstlisting}[language=Python, caption={\SemASE{} activation and round-trip persistence.}]
import semase
semase.activate(endpoint="https://kg.degu.cl/.../sparql")

from ase import Atoms
water = Atoms('H2O',
    positions=[(0, 0, 0), (0, 0.76, 0.59), (0, -0.76, 0.59)])
water.save_to_kg(uri="ase:water-001")

loaded = Atoms()
loaded.update_from_kg(uri="ase:water-001")
print(loaded.get_chemical_formula())  # H2O
\end{lstlisting}

%% ====================================================================
